%%
%% Copyright 2007, 2008, 2009 Elsevier Ltd
%%
%% This file is part of the 'Elsarticle Bundle'.
%% ---------------------------------------------
%%
%% It may be distributed under the conditions of the LaTeX Project Public
%% License, either version 1.2 of this license or (at your option) any
%% later version.  The latest version of this license is in
%%    http://www.latex-project.org/lppl.txt
%% and version 1.2 or later is part of all distributions of LaTeX
%% version 1999/12/01 or later.
%%
%% The list of all files belonging to the 'Elsarticle Bundle' is
%% given in the file `manifest.txt'.
%%

%% Template article for Elsevier's document class `elsarticle'
%% with numbered style bibliographic references
%% SP 2008/03/01

\documentclass[preprint, review, 12pt, a4paper]{elsarticle}

%% Use the option review to obtain double line spacing
%% \documentclass[authoryear,preprint,review,12pt]{elsarticle}

%% For including figures, graphicx.sty has been loaded in
%% elsarticle.cls. If you prefer to use the old commands
%% please give \usepackage{epsfig}

%% The amssymb package provides various useful mathematical symbols
\usepackage{amssymb}
%% The amsthm package provides extended theorem environments
%% \usepackage{amsthm}

%% The lineno packages adds line numbers. Start line numbering with
%% \begin{linenumbers}, end it with \end{linenumbers}. Or switch it on
%% for the whole article with \linenumbers.
\usepackage{lineno}
\usepackage{float}
\restylefloat{table}


\usepackage{hyperref}
\journal{SoftwareX}

\begin{document}

\begin{frontmatter}

%% Title, authors and addresses

%% use the tnoteref command within \title for footnotes;
%% use the tnotetext command for theassociated footnote;
%% use the fnref command within \author or \address for footnotes;
%% use the fntext command for theassociated footnote;
%% use the corref command within \author for corresponding author footnotes;
%% use the cortext command for theassociated footnote;
%% use the ead command for the email address,
%% and the form \ead[url] for the home page:
%% \title{Title\tnoteref{label1}}
%% \tnotetext[label1]{}
%% \author{Name\corref{cor1}\fnref{label2}}
%% \ead{email address}
%% \ead[url]{home page}
%% \fntext[label2]{}
%% \cortext[cor1]{}
%% \address{Address\fnref{label3}}
%% \fntext[label3]{}

\title{fastmat: Efficient Linear Transforms in Python}

%% use optional labels to link authors explicitly to addresses:
%% \author[label1,label2]{}
%% \address[label1]{}
%% \address[label2]{}

\author{S. Semper, C. W. Wagner}

\address{EMS Group, TU Ilmenau, some address}

\begin{abstract}
%% Text of abstract
We describe a new Python package that allows to define and work with linear transformations very intuitively and efficiently.

\end{abstract}

\begin{keyword}
%% keywords here, in the form: keyword \sep keyword
linear transforms \sep linear algebra \sep compressed sensing

%% PACS codes here, in the form: \PACS code \sep code

%% MSC codes here, in the form: \MSC code \sep code
%% or \MSC[2008] code \sep code (2000 is the default)

\end{keyword}

\end{frontmatter}

\section*{Required Metadata}\label{metadata}

\section*{Current code version}\label{current_code}

Ancillary data table required for subversion of the codebase. Kindly replace examples in right column with the correct information about your current code, and leave the left column as it is.

\begin{table}[H]
\begin{tabular}{|l|p{6.5cm}|p{6.5cm}|}
\hline
\textbf{Nr.} & \textbf{Code metadata description} & \textbf{Please fill in this column} \\
\hline
C1 & Current code version & v0.2 \\
\hline
C2 & Permanent link to code/repository used for this code version & For example: \url{https://github.com/EMS-TU-Ilmenau/fastmat} \\
\hline
C3 & Code Ocean compute capsule & \\
\hline
C4 & Legal Code License   & Apache License Version 2.0 \\
\hline
C5 & Code versioning system used & git \\
\hline
C6 & Software code languages, tools, and services used & Python, Cython \\
\hline
C7 & Compilation requirements, operating environments \& dependencies & \\
\hline
C8 & If available Link to developer documentation/manual & \url{https://fastmat.readthedocs.io/en/latest/} \\
\hline
C9 & Support email for questions & \url{sebastian.semper@tu-ilmenau.de}\\
\hline
\end{tabular}
\caption{Code metadata (mandatory)}
\label{mand_metadata}
\end{table}


\linenumbers


\section{Motivation and significance}\label{motivation}
% !TEX root = ../fastmat.tex

% Introduce the scientific background and the motivation for developing the software.
%
% Explain why the software is important, and describe the exact (scientific) problem(s) it solves.
%
% Indicate in what way the software has contributed (or how it will contribute in the future) to the process of scientific discovery; if available, this is to be supported by citing a research paper using the software.
%
% Provide a description of the experimental setting (how does the user use the software?).
%
% Introduce related work in literature (cite or list algorithms used, other software etc.).

\begin{itemize}
\item linear mappings are present in applied math, physics, engineering, signal processing and data science
\item represented as matrices when acting between finite dimensional spaces
\item many mappings are not arbitrary but follow a certain regularity, logic or physics: convolution, differentiation, integration
\item they are structured: fewer degrees of freedom
\item naturally: this structure is also present in our minds, our notation, our thinking
\item so the represenation as a matrix becomes wasteful and cumbersome
\item also: applying the mapping is not a matrix-vector product anymore, but an algorithm on its own, specifically tailored to the mapping
\item we have to drop the represenation as a matrix
\item we want to combine the rectangle of numbers and the specifically tailored algorithm into one concept
\item many algorithms only need the mapping and not the whole matrix (cg, bicgstab, lancosz iteration, ista, omp)
\end{itemize}




\section{Software description}\label{description}
% !TEX root = ../fastmat.tex

% Describe the software in as much as is necessary to establish a vocabulary needed to explain its impact.

\subsection{Software Architecture}\label{architecture}

% Give a short overview of the overall software architecture; provide a pictorial component overview or similar (if possible). If necessary provide implementation details.

\begin{itemize}
\item cython, numpy, scipy
\item zoo of matrices
  \begin{itemize}
    \item derived from (unstructured) base class
    \item combination to create new classes: partial(product(fourier.H, diag, fourier)) $\rightarrow$ toeplitz/circulant
  \end{itemize}
\item zoo of algorithms: class based, logging, inspection
\item cmath for high performance
\item inspecting: testing and benchmarking
\item conversion to scipy linear operator $\rightarrow$ unlock their algorithms
\end{itemize}

\subsection{Software Functionalities}\label{features}

% Present the major functionalities of the software.
\begin{itemize}
  \item zoo of matrices
  \begin{itemize}
    \item derived from (unstructured) base class
    \item interoperability between all derived classes, structure is preserved and known
    \item combination to create new classes: partial(product(fourier.H, diag, fourier)) $\rightarrow$ toeplitz/circulant
  \end{itemize}
  \item zoo of algorithms: sparse recovery
  \item conversion to scipy linear operator $\rightarrow$ unlock their algorithms
\end{itemize}

\subsection{Sample code snippets analysis (optional)}\label{snippets}



\section{Illustrative Examples}\label{examples}
% !TEX root = ../fastmat.tex

% Provide at least one illustrative example to demonstrate the major functions.
%

\begin{itemize}
  \item Ultrasonic NDT
  \begin{itemize}
    \item goal: detect defects non destructively (production, medicine, maintenance)
    \item move around an ultrasound transducer to create a synthetic aperture
    \item discrete data model: matrix(kronecker, fourier, toeplitz, partial, diagonal)
    \item sparsity: few defects, positions non zeros in vector encode defect position
    \item origin: kronecker: multidimensional data, fourier: compression, toeplitz: geometry of sampling and physics
    \item write down sparse recovery problem
    \item OMP: model dependent calculation $c = A^H r$, time/memory complexity depends on complexity of $A$, rest of linear time/memory-complexity
  \end{itemize}
  \item Goldcodes
  \begin{itemize}
    \item ...
  \end{itemize}
\end{itemize}


\section{Impact}\label{impact}
% !TEX root = ../fastmat.tex

% \textbf{This is the main section of the article and the reviewers weight the description here appropriately}
%
% Indicate in what way new research questions can be pursued as a result of the software (if any).
%
% Indicate in what way, and to what extent, the pursuit of existing research questions is improved (if so).
%
% Indicate in what way the software has changed the daily practice of its users (if so).
%
% Indicate how widespread the use of the software is within and outside the intended user group.
%
% Indicate in what way the software is used in commercial settings and/or how it led to the creation of spin-off companies (if so).



\section{Conclusions}\label{conclusion}
% !TEX root = ../fastmat.tex

% Set out the conclusion of this original software publication.



\section{Conflict of Interest}
We wish to confirm that there are no known conflicts of interest associated with this publication and there has been no significant financial support for this work that could have influenced its outcome.


\section*{Acknowledgements}\label{ack}

% Optionally thank people and institutes you need to acknowledge.

%% The Appendices part is started with the command \appendix;
%% appendix sections are then done as normal sections
%% \appendix

%% \section{}
%% \label{}

%% References:
%% If you have bibdatabase file and want bibtex to generate the
%% bibitems, please use
%%
% Please add the reference to the software repository if DOI for software  is available.

\bibliographystyle{elsarticle-num}
\bibliography{bib.bib}


\section*{Current executable software version}\label{current}

Ancillary data table required for sub version of the executable software: (x.1, x.2 etc.) kindly replace examples in right column with the correct information about your executables, and leave the left column as it is.

\begin{table}[!h]
\begin{tabular}{|l|p{6.5cm}|p{6.5cm}|}
\hline
\textbf{Nr.} & \textbf{(Executable) software metadata description} & \textbf{Please fill in this column} \\
\hline
S1 & Current software version & For example 1.1, 2.4 etc. \\
\hline
S2 & Permanent link to executables of this version  & For example: $https://github.com/combogenomics/$ $DuctApe/releases/tag/DuctApe-0.16.4$ \\
\hline
S3 & Legal Software License & List one of the approved licenses \\
\hline
S4 & Computing platforms/Operating Systems & Linux, OS X, Microsoft Windows \\
\hline
S5 & Installation requirements \& dependencies & Python, Numpy, Scipy, matplotlib\\
\hline
S6 & If available, link to user manual - if formally published include a reference to the publication in the reference list & \url{https://fastmat.readthedocs.io/en/latest/} \\
\hline
S7 & Support email for questions & \url{sebastian.semper@tu-ilmenau.de}\\
\hline
\end{tabular}
\caption{Software metadata (optional)}\label{meta_optional}
\end{table}

\end{document}
\endinput
%%
%% End of file `SoftwareX_article_template.tex'.
